\documentclass[12pt,a4paper]{article}

\usepackage[margin=1in]{geometry}

\usepackage{amsmath,amssymb,amsthm}
\usepackage{mathtools}
\usepackage{booktabs}
\usepackage{enumitem}
\usepackage{hyperref}

% Readability: open up the text
\setlength{\parskip}{0.4em plus 0.1em minus 0.1em}
\renewcommand{\arraystretch}{1.2}
\setlist{itemsep=0.3em, parsep=0.1em}

\newtheorem{theorem}{Theorem}
\newtheorem{lemma}{Lemma}
\newtheorem{corollary}{Corollary}
\newtheorem{remark}{Remark}

\newcommand{\R}{\mathbb{R}}
\newcommand{\C}{\mathbb{C}}
\newcommand{\Z}{\mathbb{Z}}
\newcommand{\Q}{\mathbb{Q}}
\newcommand{\KS}{\mathrm{KS}}
\newcommand{\Norm}[1]{N_{K/\Q}\!\left(#1\right)}
\newcommand{\Tr}[1]{\mathrm{Tr}_{K/\Q}\!\left(#1\right)}
\newcommand{\Gal}{\mathrm{Gal}}
\newcommand{\braket}[2]{\langle #1 | #2 \rangle}

\begin{document}

\title{Galois-theoretic classification of Kochen--Specker sets
in dimension three}

\author{Michael Kernaghan\\
\small Pacific Quantum Systems, Vancouver, Canada}

\date{April 2026}

\maketitle

\begin{abstract}
\noindent
We classify the two-element coordinate alphabets
$\mathcal{A} = \{0, \pm 1, \pm x\}$ with~$x$ an algebraic integer in a
quadratic extension $K/\Q$ (or $x = \pm 2$, the rational case) for
which the \emph{raw} ray set $S(\mathcal{A}) \subset \C^3$, before
cross-product completion, is KS-uncolorable.  The classification gives
three alphabet-invariant conditions, each controlled by the Galois
norm and trace of~$x$:
\begin{enumerate}
\item[(R)]  $x = \pm 2$ \emph{(rational)},
\item[(i)]  $\Tr{x} = 0$ and $|x|^2 = 2$ \emph{(pure modulus-2)},
\item[(ii)] $\Norm{x} = 1$ and $|\Tr{x}| = 1$ \emph{(phase)}.
\end{enumerate}
Sufficiency is witnessed by explicit KS constructions (CK-31 for~(R),
Peres-33 and the $\Z[\sqrt{-2}]$-33 set for~(i), the Eisenstein-33 set
for~(ii)).  Necessity is proved by a finite enumeration of three-term
vanishing sums over the product set $\bar{\mathcal{A}} \cdot
\mathcal{A}$, combined with SMT verification of the resulting
colorability obstruction, showing that when no condition holds the
orthogonality graph admits an explicit $\{0,1\}$-coloring.  As a
corollary we obtain a norm-and-trace characterization of the raw
KS-contextual quadratic alphabets.  A parallel survey under cross-product
completion~\cite{Kernaghan2026islands} identifies three further
algebraic islands (Gaussian, Heegner-7, golden ratio) that fall outside
the raw classification; establishing a Galois-theoretic classification
of the completed six-island picture remains an open problem.  The
classification connects to the result of Cortez, Morales, and Reyes
that $\Z[1/6]$ is the minimal ring supporting KS contextuality in
dimension~3.
\end{abstract}

\section{Introduction}

The Kochen--Specker (KS) theorem~\cite{KochenSpecker1967} establishes
that quantum measurements in Hilbert spaces of dimension $d \geq 3$
cannot be explained by noncontextual hidden-variable theories.  The proof
proceeds by exhibiting a finite set of rays---a \emph{KS set}---that
admits no consistent $\{0,1\}$-valued coloring respecting the rules: (I)
orthogonal rays cannot both receive the value~1, and (II) every
orthogonal triple (triad) contains exactly one~1.  The existence of such
a set refutes the assumption that measurement outcomes pre-exist
independently of which other compatible measurements are performed
alongside them.

The starting observation behind the present work is visible in the
cube diagrams of Peres's textbook~\cite{Peres1993}: two of the three
known three-dimensional KS constructions---the 31-vector
Conway--Kochen set (CK-31) and the 33-vector Peres set---are displayed
side by side, yet their coordinate arithmetics differ in a suggestive
way.  CK-31 uses the integer alphabet $\{0, \pm 1, \pm 2\}$, with
every vertex labelled by $1$s and $2$s.  The Peres set uses
$\{0, \pm 1, \pm\sqrt{2}\}$, replacing the ``$2$'' with ``$\sqrt{2}$.''
Both share the same two-element template $\{0, \pm 1, \pm x\}$; only
the generator~$x$ changes.  This pattern invites a natural question:
\emph{what happens if we vary~$x$ systematically?}  Could coordinates
built from $1/\!\sqrt{2}$, or $\sqrt{3}$, or a cube root of~$2$, or a
sixth root of unity produce new---perhaps smaller---KS sets?

A computational survey~\cite{Kernaghan2026islands} pursued this
question across quadratic, cyclotomic, and golden-ratio number fields,
fixing a two-element coordinate alphabet $\mathcal{A} = \{0, \pm 1,
\pm x\}$ with~$x$ a ring-of-integers generator, generating the
induced ray set, and testing for KS-uncolorable subsets.  The survey
found six distinct algebraic \emph{islands}:
\begin{itemize}
  \item $\Z$: integer alphabet, CK-31 (31 vectors);
  \item $\Z[\sqrt{2}]$: Peres-type set (33 vectors);
  \item $\Z[\omega]$, $\omega = e^{2\pi i/3}$ (Eisenstein integers):
        Cabello-type set (33 vectors);
  \item $\Z[\sqrt{-2}]$: 33 vectors;
  \item $\Z\!\left[\tfrac{1+\sqrt{-7}}{2}\right]$ (Heegner-7 ring): 43
        vectors;
  \item $\Q(\varphi)$, $\varphi = \tfrac{1+\sqrt{5}}{2}$ (golden-ratio
        field): 52 vectors.
\end{itemize}
The survey identified two recurring mechanisms driving KS-uncolorability:
a \emph{modulus-2} mechanism ($|x|^2 = 2$ in the first three cases above)
and a \emph{phase} mechanism ($1 + \omega + \omega^2 = 0$ in the
Eisenstein case, a degenerate instance of a unit vanishing sum).  These
patterns were, however, established computationally and without a
unifying algebraic explanation.

This paper closes that gap for the family of quadratic extensions.  We
prove that the two mechanisms are the \emph{only} ones capable of
producing KS-uncolorable \emph{raw} ray sets from a two-element
quadratic alphabet,
and that each is characterized precisely by the Galois norm and trace of
the ring generator.  The proof classifies
three-term vanishing sums over the product alphabet
$\bar{\mathcal{A}} \cdot \mathcal{A}$, determines the resulting
orthogonality-graph structure, and shows that when neither mechanism is
active the constraint graph decomposes into mutually independent
subproblems each admitting an explicit valid coloring.

\begin{theorem}\label{thm:main}
Let $\mathcal{A} = \{0,\pm 1,\pm x\}$ with $x \notin \{0, \pm 1\}$ an
algebraic integer in a quadratic extension $K = \Q(\sqrt d)/\Q$, or
$x = \pm 2$ (the degenerate rational case $K = \Q$).
The \emph{raw} ray set $S(\mathcal{A}) \subset \C^3$ is KS-uncolorable
if and only if one of:
\begin{enumerate}
  \item[(R)]  $x = \pm 2$ \emph{(rational degenerate case)};
  \item[(i)]  $\Tr{x} = 0$ and $|x|^2 = 2$ \emph{(pure modulus-2:
    $x = \pm\sqrt{2}$ or $x = \pm\sqrt{-2}$)};
  \item[(ii)] $\Norm{x} = 1$ and $|\Tr{x}| = 1$ \emph{(phase, i.e.\
    $1 + x + \bar x = 0$ or $1 - x - \bar x = 0$)}.
\end{enumerate}
Here $\Norm{x}$ and $\Tr{x}$ denote the Galois norm and trace of~$x$
(Section~\ref{sec:setup}).  The three conditions are each invariant under
$x \mapsto -x$ and under $x \mapsto \bar x$, hence depend only on the
alphabet~$\mathcal{A}$ itself.
\end{theorem}

\begin{remark}[Raw vs.\ completed]\label{rem:raw-scope}
Theorem~\ref{thm:main} classifies the \emph{raw} ray set
$S(\mathcal{A})$, before cross-product completion.  A systematic
computational survey~\cite{Kernaghan2026islands} identifies
\emph{six} algebraic islands in dimension three once the ray pool is
closed under cross products: the three captured above, together with
the Gaussian alphabet ($x = 1+i$), the Heegner-7 alphabet
($x = (1+\sqrt{-7})/2$), and the golden-ratio alphabet ($x = \varphi$).
The latter three are \emph{colorable} at the raw level (insufficient
triad density) and become KS-uncolorable only after completion
introduces additional rays.  Classifying the completed six-island
picture by Galois-theoretic invariants remains an open problem.
\end{remark}

The proof occupies Section~\ref{sec:proof}.  Section~\ref{sec:consequences}
derives corollaries including the imaginary quadratic classification
and the Galois symmetry of the orthogonality graph.  Section~\ref{sec:discussion} discusses
connections to the Cortez--Morales--Reyes ring-theoretic result.

\section{Setup}\label{sec:setup}

\subsection{Quadratic extensions and their arithmetic}

Let $d \in \Z$ be squarefree and $K = \Q(\sqrt{d})$ the associated
quadratic extension.  The ring of integers $\mathcal{O}_K$ has the form
$\Z[x]$ where the generator~$x$ is:
\[
  x = \begin{cases}
    \sqrt{d} & \text{if } d \equiv 2, 3 \pmod{4}, \\
    \dfrac{1 + \sqrt{d}}{2} & \text{if } d \equiv 1 \pmod{4}.
  \end{cases}
\]
The Galois group $\Gal(K/\Q) = \{1, \sigma\}$ is generated by the
non-trivial automorphism $\sigma(\sqrt{d}) = -\sqrt{d}$.  From this we
define the \emph{Galois norm} and \emph{Galois trace} of any $\alpha \in
K$:
\[
  \Norm{\alpha} = \alpha \cdot \sigma(\alpha), \qquad
  \Tr{\alpha} = \alpha + \sigma(\alpha).
\]
For the generator~$x$ these take concrete values depending on~$d$.  In
the case $d \equiv 2, 3 \pmod{4}$: $\Norm{x} = \sqrt{d} \cdot
(-\sqrt{d}) = -d$ and $\Tr{x} = 0$.  In the case $d \equiv 1 \pmod{4}$:
$\Norm{x} = \tfrac{1+\sqrt{d}}{2} \cdot \tfrac{1-\sqrt{d}}{2} =
\tfrac{1-d}{4}$ and $\Tr{x} = 1$.

\subsection{The coordinate alphabet and ray set}

Fix a quadratic extension $K/\Q$ with generator~$x$ as above.  Define
the \emph{two-element coordinate alphabet}
\[
  \mathcal{A} = \{0,\; +1,\; -1,\; +x,\; -x\}.
\]
The \emph{ray set} $S(\mathcal{A})$ consists of all projectively distinct
nonzero vectors $v \in \mathcal{A}^3 \setminus \{0\}$, where two vectors
$v, w$ represent the same ray if $v = \lambda w$ for some scalar
$\lambda \in \C^\times$.  Orthogonality is defined via the standard
Hermitian inner product:
\[
  \braket{v}{w} = \sum_{k=1}^{3} \bar{v}_k w_k = 0.
\]

\subsection{The product set and vanishing sums}

For the orthogonality analysis the key object is the \emph{product set}
\[
  \bar{\mathcal{A}} \cdot \mathcal{A}
  = \{\bar{a} b : a, b \in \mathcal{A}\}
  = \{0,\; \pm 1,\; \pm x,\; \pm \bar{x},\; \pm |x|^2\}.
\]
The inner product $\braket{v}{w}$ for $v, w \in \mathcal{A}^3$ is a sum
of at most three terms, each drawn from $\bar{\mathcal{A}} \cdot
\mathcal{A}$.  Orthogonality $\braket{v}{w} = 0$ therefore reduces to a
\emph{vanishing sum} of one, two, or three elements from this product
set.

A \emph{primitive three-term vanishing sum} from
$\bar{\mathcal{A}} \cdot \mathcal{A}$ is an identity
$\alpha + \beta + \gamma = 0$ with $\alpha, \beta, \gamma \in
\bar{\mathcal{A}} \cdot \mathcal{A}$, not reducible to a shorter
vanishing sum.  The classification of such sums---there are finitely many
modulo sign and permutation---drives the proof of
Theorem~\ref{thm:main}.

Explicitly, the product set elements are controlled by two invariants of
the extension:
For imaginary quadratic fields, complex conjugation coincides with
$\sigma$ on~$K$, so $|x|^2 = x\bar{x} = |\Norm{x}|$ and $x + \bar{x}
= \Tr{x}$; for real quadratic fields, $\bar{x} = x$ and $|x|^2 = x^2$.
Condition~(i) $|x|^2 = 2$ means the product set contains an element of
absolute value~$2$, enabling cancellations $1 + 1 - |x|^2 = 0$.
Condition~(ii) means $1 + x + \bar{x} = 1 + \Tr{x} = 0$, a
unit-modulus vanishing sum as in $1 + \omega + \omega^2 = 0$.

\section{Main theorem}\label{sec:proof}

We restate the main result for convenience.

\begin{theorem}[Galois classification, restated]\label{thm:galois}
Let $\mathcal{A} = \{0,\pm 1,\pm x\}$ with $x \notin \{0,\pm 1\}$ an
algebraic integer in a quadratic extension $K = \Q(\sqrt d)/\Q$, or
$x = \pm 2$ (rational case).  The raw ray set
$S(\mathcal{A}) \subset \C^3$ is KS-uncolorable if and only if one of:
\begin{enumerate}
  \item[\emph{(R)}]  $x = \pm 2$ \emph{(rational)};
  \item[\emph{(i)}]  $\Tr{x} = 0$ and $|x|^2 = 2$ \emph{(pure modulus-2)};
  \item[\emph{(ii)}] $\Norm{x} = 1$ and $|\Tr{x}| = 1$ \emph{(phase)}.
\end{enumerate}
\end{theorem}

\noindent\textit{Sufficiency.}
Each condition is witnessed by an explicit KS construction.
\begin{itemize}
\item Condition~(R) $x = \pm 2$ gives the 31-vector Conway--Kochen set
  CK-31 (rigidity proved by
  Trandafir--Cabello~\cite{TrandafirCabello2025}).
\item Condition~(i) is realized uniquely (up to sign) by
  $x = \pm\sqrt{2}$ and $x = \pm\sqrt{-2}$.  The real case gives the
  33-vector Peres set~\cite{Peres1991}; the imaginary case gives the
  33-vector $\Z[\sqrt{-2}]$ set~\cite{Kernaghan2026islands}, which is
  graph-isomorphic to Peres.
\item Condition~(ii) is realized by $x = \omega = e^{2\pi i/3} =
  (-1{+}\sqrt{-3})/2$ ($\Norm{\omega} = 1$, $\Tr{\omega} = -1$),
  yielding the 33-vector Eisenstein set~\cite{Cabello2025simplest};
  equivalently, by the canonical generator
  $x = (1{+}\sqrt{-3})/2$ ($\Norm{x} = 1$, $\Tr{x} = +1$), which
  produces the same alphabet $\{0, \pm 1, \pm x\} = \{0, \pm 1, \pm
  \bar\omega\}$.
\end{itemize}
Other alphabets with $|x|^2 = 2$ but $\Tr{x} \neq 0$ (the Gaussian
generator $x = 1+i$ and the Heegner-7 generator
$x = (1{+}\sqrt{-7})/2$) produce ray pools with orthogonal pairs but
\emph{insufficient triad density} to force uncolorability at the raw
level (Lemma~\ref{lem:sparsity}, Step~1).  Known
KS constructions over these
alphabets~\cite{Kernaghan2026islands} require cross-product completion
and lie outside the scope of the present raw classification; see
Remark~\ref{rem:raw-scope}.  \hfill$\square$

\medskip

\noindent\textit{Necessity.}
We prove the contrapositive: if neither~(i) nor~(ii) holds, then
$S(\mathcal{A})$ admits a valid $\{0,1\}$-coloring.  The argument
proceeds through two lemmas.

\begin{lemma}[Vanishing sum enumeration]\label{lem:vanishing}
Let $x \notin \{0,\pm 1\}$ be a ring-of-integers generator with
conjugate $\bar{x} = \sigma(x)$.  The primitive three-term vanishing
sums $t_1 + t_2 + t_3 = 0$ with each $t_k \in
\{0,\pm 1,\pm x,\pm \bar{x},\pm |x|^2\}$ are, up to sign and
permutation, exactly those listed in Table~\ref{tab:vanishing}.  In
particular, every such sum constrains $(N,T) = (\Norm{x},\Tr{x})$ to
one of the rows in the table.
\end{lemma}

\begin{table}[h]
\centering
\caption{Primitive three-term vanishing sums from the product set
$\bar{\mathcal{A}} \cdot \mathcal{A}$.  Here
$N = \Norm{x}$ and $T = \Tr{x}$.}\label{tab:vanishing}
\begin{tabular}{@{}lll@{}}
\toprule
Pattern & Constraint & $N,\; T$ \\
\midrule
$1 + 1 + (-x) = 0$ & $x = 2$ & (rational) \\
$1 + 1 + (-|x|^2) = 0$ & $|x|^2 = 2$ & $\pm 2$, varies \\
$1 + x + \bar{x} = 0$ & $\Tr{x} = -1$ & varies,\; ${-1}$ \\
$1 - x - \bar{x} = 0$ & $\Tr{x} = 1$ & varies,\; $1$ \\
$1 + x + (-|x|^2) = 0$ & $|x|^2 = 1 + x$ & ${-1},\; 1$ (golden) \\
$|x|^2 + |x|^2 + (-1) = 0$ & $|x|^2 = \tfrac{1}{2}$ & $\equiv$ row~2 \\
$x + x + (-|x|^2) = 0$ & $|x|^2 = 2x$ & $\equiv$ row~1 \\
\bottomrule
\end{tabular}
\end{table}

\begin{proof}
The product set $\bar{\mathcal{A}} \cdot \mathcal{A}$ consists (up to
sign) of $\{0, 1, x, \bar{x}, |x|^2\}$.  A primitive three-term
vanishing sum $\alpha + \beta + \gamma = 0$ with each term drawn
from $\pm\{1, x, \bar{x}, |x|^2\}$ must satisfy two independent
constraints: the \emph{magnitude} constraint (the three terms must
form a closed triangle in~$\C$) and the \emph{algebraic} constraint
(the sum must vanish exactly in the ring $\Z[x, \bar{x}]$).

Writing each term as $\pm s_k$ with $s_k \in \{1, x, \bar{x},
|x|^2\}$, the vanishing sum becomes a polynomial identity in~$x$
and~$\bar{x}$.  Since $\{1, x\}$ is a $\Q$-basis of~$K$ and
$\bar{x} = T - x$ (where $T = \Tr{x}$), every product-set element
has a unique representation $a + bx$ with $a, b \in \Q$.  The
representations depend on whether~$K$ is real or imaginary quadratic.

\smallskip
\noindent\emph{Imaginary case} ($d < 0$): complex conjugation
coincides with~$\sigma$ on~$K$, so $|x|^2 = x\bar{x} = \Norm{x}$:
\[
  1 \mapsto (1, 0),\quad x \mapsto (0, 1),\quad
  \bar{x} \mapsto (T, -1),\quad |x|^2 \mapsto (N, 0).
\]

\noindent\emph{Real case} ($d > 0$): $x$ is real, so $\bar{x} = x$
in the product set (the entries $\pm\bar{x}$ collapse to $\pm x$),
and $|x|^2 = x^2 = Tx - N$ via the minimal polynomial
$x^2 - Tx + N = 0$:
\[
  1 \mapsto (1, 0),\quad x \mapsto (0, 1),\quad
  \bar{x} = x \mapsto (0, 1),\quad |x|^2 = x^2 \mapsto (-N, T).
\]

\noindent In the enumeration below, rows involving~$\bar{x}$ as
distinct from~$x$ arise only in the imaginary case.  Rows
involving~$|x|^2$ are labeled (I) for imaginary and (R) for real
where the constraints differ.

A vanishing sum $\alpha + \beta + \gamma = 0$ with each
$\alpha, \beta, \gamma \in \pm\{1, x, \bar{x}, |x|^2\}$ yields
two rational equations: the coefficient of~$1$ vanishes and the
coefficient of~$x$ vanishes.  The four unsigned elements
$\{1, x, \bar{x}, |x|^2\}$ give exactly $\binom{4+2}{3} = 20$
unordered triples with repetition.  We enumerate all~$20$, with all
sign assignments $(\varepsilon_1, \varepsilon_2,
\varepsilon_3) \in \{\pm 1\}^3$ modulo the overall sign and
permutation symmetry:

\smallskip
\begin{center}
\small
\renewcommand{\arraystretch}{1.25}
\begin{tabular}{@{}llll@{}}
\toprule
Triple & Coeff.\ of~$1$ & Coeff.\ of~$x$ & Solution \\
\midrule
\multicolumn{4}{@{}l}{\textit{Triples from $\{1\}$ only:}} \\[2pt]
$\{1, 1, 1\}$ & $\pm 1 \pm 1 \pm 1 = 0$ & $0 = 0$ & impossible \\[4pt]
\multicolumn{4}{@{}l}{\textit{Triples from $\{1, x\}$:}} \\[2pt]
$\{1, 1, x\}$ & $\pm 1 \pm 1 = 0$ & $\pm 1 = 0$ & impossible \\
$\{1, x, x\}$ & $\pm 1 = 0$ & $\pm 1 \pm 1 = 0$ & impossible \\
$\{x, x, x\}$ & $0 = 0$ & $\pm 1 \pm 1 \pm 1 = 0$ & impossible \\[4pt]
\multicolumn{4}{@{}l}{\textit{Triples involving $\bar{x}$ (imaginary only, since $\bar{x} = x$ for real):}} \\[2pt]
$\{1, 1, \bar{x}\}$ & $\pm 1 \pm 1 \pm T = 0$ & $\mp 1 = 0$ & impossible \\
$\{1, x, \bar{x}\}$ & $1 + T = 0$ & $1 - 1 = 0$ & $T = -1$ (Row~3a) \\
$\{1, x, \bar{x}\}$ & $1 - T = 0$ & $-1 + 1 = 0$ & $T = 1$ (Row~3b) \\
$\{1, \bar{x}, \bar{x}\}$ & $\varepsilon_1 + (\varepsilon_2{+}\varepsilon_3)T = 0$ & $-(\varepsilon_2{+}\varepsilon_3) = 0$ & impossible ($\varepsilon_1 = 0$) \\
$\{x, x, \bar{x}\}$ & $T = 0$ & $\pm 2 \mp 1 = 0$ & impossible \\
$\{x, \bar{x}, \bar{x}\}$ & $(\varepsilon_2{+}\varepsilon_3)T = 0$ & $\varepsilon_1 - (\varepsilon_2{+}\varepsilon_3) = 0$ & impossible$^{\ddagger}$ \\
$\{\bar{x}, \bar{x}, \bar{x}\}$ & $(\sum\varepsilon_k)T = 0$ & $-\sum\varepsilon_k = 0$ & impossible ($\sum\varepsilon_k \neq 0$) \\[4pt]
\multicolumn{4}{@{}l}{\textit{Triples involving $|x|^2$\; --- split (I)\,=\,imaginary, (R)\,=\,real:}} \\[2pt]
$\{1, 1, |x|^2\}$ (I) & $1 + 1 - N = 0$ & $0 = 0$ & $N = 2$ (Row~2) \\
$\{1, 1, |x|^2\}$ (R) & $1 + 1 + N = 0$ & $-T = 0$ & $N = {-2},\, T = 0$ (Row~2) \\
$\{1, x, |x|^2\}$ (I) & $1 - N = 0$ & $1 = 0$ & impossible \\
$\{1, x, |x|^2\}$ (R) & $1 + N = 0$ & $1 - T = 0$ & $N = {-1},\, T = 1$ (Row~4) \\
$\{1, \bar{x}, |x|^2\}$ (I) & $1 + T - N = 0$ & $-1 = 0$ & impossible \\
$\{x, x, |x|^2\}$ (I) & $\varepsilon_3 N = 0$ & $\varepsilon_1{+}\varepsilon_2 = 0$ & impossible ($N \neq 0$) \\
$\{x, x, |x|^2\}$ (R) & $-\varepsilon_3 N = 0$ & $\varepsilon_1{+}\varepsilon_2{+}\varepsilon_3 T = 0$ & impossible ($N \neq 0$) \\
$\{x, \bar{x}, |x|^2\}$ (I) & $T - N = 0$ & $0 = 0$ & $N = T$\;$^\dagger$ \\
$\{\bar{x}, \bar{x}, |x|^2\}$ (I) & $2T - N = 0$ & $-2 = 0$ & impossible \\
$\{\bar{x}, |x|^2, |x|^2\}$ (I) & $\varepsilon_1 T + (\varepsilon_2{+}\varepsilon_3) N = 0$ & $-\varepsilon_1 = 0$ & impossible \\
$\{|x|^2, |x|^2, 1\}$ (I) & $2N - 1 = 0$ & $0 = 0$ & $N = 1/2$\;($\equiv$ Row~2) \\
$\{|x|^2, |x|^2, 1\}$ (R) & $-2N + 1 = 0$ & $2T = 0$ & $N = 1/2,\, T = 0$\;($\equiv$ Row~2) \\
$\{|x|^2, |x|^2, x\}$ (I) & $2N = 0$ & $1 = 0$ & impossible \\
$\{|x|^2, |x|^2, x\}$ (R) & $-2N = 0$ & $2T + 1 = 0$ & impossible ($N \neq 0$) \\
$\{|x|^2, |x|^2, |x|^2\}$ (I) & $\pm N \pm N \pm N = 0$ & $0 = 0$ & $N = 0$ (excluded) \\
\bottomrule
\end{tabular}
\end{center}
\smallskip

\noindent $^\ddagger$For $\{x, \bar{x}, \bar{x}\}$: if $T \neq 0$,
then $\varepsilon_2 + \varepsilon_3 = 0$ forces $\varepsilon_1 = 0$;
if $T = 0$, the $x$-equation gives
$\varepsilon_1 = \varepsilon_2 + \varepsilon_3 \in \{0, \pm 2\}$,
impossible since $\varepsilon_1 = \pm 1$.

\noindent $^\dagger$The case $\{x, \bar{x}, |x|^2\}$ arises only in
the imaginary case ($\bar{x} \neq x$).  The constraint $N = T$
gives $|\Norm{x}| = \Tr{x}$: for $d \equiv 1 \pmod{4}$,
$(1{-}d)/4 = 1$ gives $d = -3$ (Row~3a via a different presentation).
The \emph{rational case} $x = 2$ ($K = \Q$) does not admit a
$\Q$-basis decomposition; it is verified directly: $1 + 1 - 2 = 0$
(Row~1).  This accounts for all $20$ unordered triples; every
surviving solution maps to one of the seven rows of
Table~\ref{tab:vanishing}, completing the enumeration.

Row~3b merits comment: the sign pattern
$(\varepsilon_1, \varepsilon_2, \varepsilon_3) = (1, -1, -1)$ on the
triple $\{1, x, \bar{x}\}$ gives $1 - x - \bar{x} = 1 - T = 0$,
hence $T = 1$ with no constraint on~$N$.  This vanishing sum
holds for all imaginary quadratic fields with
$d \equiv 1 \pmod{4}$ (e.g., $d = -11, -15, -19, \ldots$).
Unlike Row~3a ($T = -1$, which produces the Eisenstein island at
$d = -3$), Row~3b does not lead to KS-uncolorability:
the all-nonzero triads it generates are structurally
decoupled from the one-zero triads, and the full ray set remains
colorable.  This is proved for all $N \geq 3$ by SMT
verification~\cite{deMouraBjorner2008}; see
Lemma~\ref{lem:sparsity} for the proof.

Row~4 merits comment: the constraint $|x|^2 = 1 + x$ is satisfied by
$x = \varphi = \tfrac{1+\sqrt{5}}{2}$ (golden ratio), which has
$\Norm{\varphi} = -1$ and $\Tr{\varphi} = 1$.  This does not satisfy
(i) or~(ii) of Theorem~\ref{thm:galois}.  Although
$\Q(\varphi)$-alphabet KS sets exist~\cite{Kernaghan2026islands} with
52~vectors, they require cross-product completion beyond the initial
alphabet $\mathcal{A}$ and thus lie outside the scope of the present
classification.
\end{proof}

\begin{lemma}[Colorability of the complement]\label{lem:sparsity}
Suppose $x$ satisfies none of conditions~\emph{(R)},~\emph{(i)},
or~\emph{(ii)} of Theorem~\ref{thm:galois}.  Then the ray set
$S(\mathcal{A})$ admits an explicit valid $\{0,1\}$-coloring.
\end{lemma}

\noindent\textit{Strategy.}
The proof splits the residual case into two subfamilies according to
whether the Row-2 vanishing sum $1 + 1 - |x|^2 = 0$ is available.

\smallskip
\noindent\textbf{Subfamily A} (\emph{no Row-2 cancellation}):
$|x|^2 \neq 2$.  Then by Lemma~\ref{lem:vanishing} the only surviving
vanishing sum is Row~3b ($\Tr{x} = 1$, available when $d \equiv 1
\pmod{4}$ imaginary), and Row~4 (golden ratio, $x = \varphi$,
$|x|^2 = 1 + x$) which falls outside~$\mathcal{O}_K$ for any
$d \neq 5$ and is treated separately in Section~\ref{sec:discussion}.
We show that the all-nonzero triads from Row~3b are structurally
decoupled from the one-zero triads, and no type-B triad (mixing the
two families) exists.  The decoupling reduces colorability to two
independent subproblems, both of which are explicitly colorable.

\smallskip
\noindent\textbf{Subfamily B} (\emph{Row-2 cancellation, but
$\Tr{x} \neq 0$}): $|x|^2 = 2$ and $\Tr{x} \neq 0$.  This is the
\emph{non-pure} modulus-2 case, where condition~(i) of
Theorem~\ref{thm:galois} fails (it requires $\Tr{x} = 0$ as well).
Solving $|a + b x|^2 = 2$ over $\mathcal{O}_K$ with $\Tr{x} \neq 0$,
the alphabet is---up to sign and Galois-conjugation---one of exactly
two:
\begin{itemize}
  \item the Gaussian alphabet $\{0,\pm 1,\pm(1+i)\}$ over
    $K = \Q(i)$ (standard generator, $\Tr{x} = 2$);
  \item the Heegner-7 halfint alphabet $\{0,\pm 1,\pm(1+\sqrt{-7})/2\}$
    over $K = \Q(\sqrt{-7})$ ($\Tr{x} = 1$).
\end{itemize}
For each of these two alphabets we exhibit an explicit valid
$\{0,1\}$-coloring (Table~\ref{tab:explicit-colorings}), verified
both by direct constraint check and by an independent SAT solver
(Glucose4, PySAT~\cite{AudemardSimon2018}).  These two cases
exhaust subfamily~B; no further analytical argument is needed.
The coexistence of type-B triads (which exist here, by the Row-2
sum) with valid colorings shows that Row-2 cancellation alone is
not a sufficient KS-uncolorability mechanism; the additional
constraint $\Tr{x} = 0$---as in condition~(i)---is required to
pack enough triads to force unsatisfiability.

\begin{table}[h]
\centering
\caption{Explicit valid $\{0,1\}$-colorings (subfamily~B).
Each row lists the green rays ($f(v) = 1$) of a coloring of the
49-ray pool of the corresponding alphabet (the remaining rays are
red, $f(v) = 0$).  Both colorings have been verified by direct
constraint check (78~Hermitian-orthogonal pair constraints,
14~triad-completeness constraints) and independently by SAT
(Glucose4, PySAT~\cite{AudemardSimon2018}); reproducer scripts are
included with this paper.}\label{tab:explicit-colorings}
\small
\begin{tabular}{@{}p{0.18\textwidth}p{0.74\textwidth}@{}}
\toprule
Alphabet & Green rays (representative vectors in $\mathcal A^3$) \\
\midrule
Gaussian \newline $x = 1+i$, $K=\Q(i)$ &
  $(1,0,0)$,
  $(1,0,1)$,
  $(1,1,0)$,
  $(1,1,1{+}i)$,
  $(1,1{+}i,1)$,
  $(1,1{+}i,1{+}i)$,
  $(1,1{+}i,-1{-}i)$,
  $(1,-1{-}i,1{+}i)$,
  $(1,-1{-}i,-1{-}i)$,
  $(1{+}i,1,1)$,
  $(1{+}i,1,-1)$
  \\
\midrule
Heegner-7 halfint \newline $x = (1{+}\sqrt{-7})/2$, \newline
$K = \Q(\sqrt{-7})$ &
  $(0,0,1)$,
  $(0,1,1)$,
  $(1,0,1)$,
  $(1,1,x)$,
  $(1,-1,-x)$,
  $(1,-x,1)$,
  $(x,1,1)$
  \\
\bottomrule
\end{tabular}
\end{table}

\begin{proof}[Proof of Lemma~\ref{lem:sparsity}]
\noindent\textit{Subfamily~B} ($|x|^2 = 2$, $\Tr{x} \neq 0$).
By the count above, the alphabet is one of two specific cases.
Table~\ref{tab:explicit-colorings} exhibits an explicit valid
$\{0,1\}$-coloring for each, completing this subfamily.  We turn to
subfamily~A.

\smallskip
\noindent\textit{Subfamily~A} ($|x|^2 \neq 2$).
By Lemma~\ref{lem:vanishing}, Rows~1--3a and~5--6 are excluded by
the negation of~(R), (i), and~(ii).  Row~4 ($x = \varphi$) is
treated in Section~\ref{sec:discussion}.  For the remaining
generators, the only surviving vanishing sum is Row~3b:
$1 - x - \bar{x} = 0$, which holds precisely when $\Tr{x} = 1$
(imaginary quadratic, $d \equiv 1 \pmod{4}$).  When $\Tr{x} \neq 1$,
no vanishing sum exists at all.  The remaining proof proceeds in
three steps.

\medskip
\noindent\textit{Step~1 (All-nonzero triads are harmless).}
When $\Tr{x} \neq 1$ (including all real quadratic and all imaginary
quadratic with $d \equiv 2, 3 \pmod{4}$), no primitive three-term
vanishing sum from $\bar{\mathcal{A}} \cdot \mathcal{A}$ exists,
so no two all-nonzero rays in $S(\mathcal{A})$ are orthogonal.

When $\Tr{x} = 1$, the vanishing sum $1 - x - \bar{x} = 0$ does
produce all-nonzero orthogonal pairs and, indeed, complete
all-nonzero triads.  Concretely, the triad
$\{(1, 1, x),\; (1, -x, -1),\; (x, -1, 1)\}$ is mutually orthogonal
(each inner product reduces to $1 - x - \bar{x} = 0$).
However, these triads involve \emph{only} all-nonzero rays and share
no rays with one-zero triads.  That no type-B triad (mixing
all-nonzero and one-zero rays) exists is proved universally: the
Hermitian cross product of any orthogonal (all-nonzero, one-zero) pair
produces a vector whose $\Q$-basis components cannot simultaneously
lie in~$\mathcal{A}$ for any integer $N \geq 3$, certified by the Z3
SMT solver~\cite{deMouraBjorner2008} over all 1024 entry patterns and
all integer $(N, T)$ with $N \geq 3$.

Since the all-nonzero and one-zero families share no rays,
a $\{0,1\}$-coloring is valid if and only if it is valid on each
family separately.  Both families are independently colorable
(Step~3 below gives 192~colorings for the one-zero family; the
all-nonzero triads, being pairwise disjoint, are trivially colorable).
The all-nonzero triads are therefore \emph{harmless}---they do not
contribute to KS-uncolorability.

\medskip
\noindent\textit{Step~2 (No type-B triads).}
Call a triad \emph{type-B} if it contains at least one all-nonzero ray
alongside rays with a zero coordinate.  An all-nonzero ray
$v \in \{\pm 1, \pm x\}^3$ can be orthogonal to a one-zero ray~$w$
(say $w_3 = 0$) via two-term cancellation:
$\bar{v}_1 w_1 + \bar{v}_2 w_2 = 0$.  However, the third member of
any triad containing $v$ and $w$ is the unique ray in the Hermitian
orthogonal complement of $\mathrm{span}(v, w)$.  In $\C^3$, this ray
is represented by the conjugate of the standard cross product,
\[
  u_k = \overline{v_{k+1} w_{k+2} - v_{k+2} w_{k+1}}
      = \bar v_{k+1}\bar w_{k+2} - \bar v_{k+2}\bar w_{k+1}
  \quad (\text{indices mod } 3),
\]
which satisfies $\braket{u}{v} = \braket{u}{w} = 0$ for the
Hermitian product $\braket{a}{b} = \sum_k \bar a_k b_k$ because
$\braket{u}{v} = \sum_k (v \times w)_k v_k = \det[v, v, w] = 0$ and
similarly $\braket{u}{w} = \det[v, w, w] = 0$.  (For real fields,
$\bar{v}_k = v_k$ and $u$ reduces to the standard bilinear cross
product.)  We show that $u \notin \mathcal{A}^3$ up to projective
equivalence.

With $w_3 = 0$, the (conjugated) cross product gives
$u = (-\bar{v}_3 \bar w_2,\; \bar{v}_3 \bar w_1,\;
\bar{v}_1 \bar w_2 - \bar{v}_2 \bar w_1)$.  Writing $\bar w_k
\in \{\pm 1, \pm \bar x\}$ and $\bar v_k \in \{\pm 1, \pm\bar x\}$,
the first two components lie in $\bar{\mathcal{A}} \cdot
\bar{\mathcal{A}}$ and the third is a difference of two such
products.  For real quadratic ($\bar x = x$) this recovers
$(-v_3 w_2, v_3 w_1, v_1 w_2 - v_2 w_1)$ as before.  For imaginary
quadratic with $d \equiv 2, 3 \pmod 4$ (standard generator,
$\bar x = -x$), the conjugate alphabet satisfies
$\bar{\mathcal A} = \mathcal A$ as a set and the analysis proceeds
by direct substitution.  For $d \equiv 1 \pmod 4$ (canonical
halfint generator, $\bar x = 1 - x \neq \pm x$), the components
of~$u$ lie a priori in $\bar{\mathcal A}^3$, and the test whether
$[u]$ is projectively equivalent to a ray in~$S(\mathcal{A})$
requires comparing $u$ to some $\lambda w'$ with
$w' \in \mathcal{A}^3$ and $\lambda \in K^\times$; since $K$ is
closed under conjugation we may equivalently test the conjugate
configuration $\bar u = v \times w$ against
$\bar\lambda \bar{\mathcal A}^3$, which is the original cross-product
analysis over the conjugate alphabet.  We will present the
calculation uniformly, flagging the $d \equiv 1 \pmod 4$ specializations
where they diverge.

\emph{Real case} ($d > 0$, so $\bar{v}_k = v_k$).
The cross-product components are $u_1 = -v_3 w_2$,
$u_2 = v_3 w_1$, $u_3 = v_1 w_2 - v_2 w_1$, with
$v_k \in \{\pm 1, \pm x\}$ and $w_k \in \{\pm 1, \pm x\}$.
For $u \in \lambda\mathcal{A}^3$, \emph{all three} components must
simultaneously lie in $\{0, \pm\lambda, \pm\lambda x\}$.

The first two components determine~$\lambda$.  Since $u_1 = -v_3 w_2$
and $u_2 = v_3 w_1$ are both nonzero (products of nonzero alphabet
entries), each is $\pm\lambda$ or $\pm\lambda x$, giving four cases
for the pair $(u_1, u_2)$:
\begin{center}
\renewcommand{\arraystretch}{1.1}
\begin{tabular}{@{}lll@{}}
$u_1 = \pm\lambda$, $u_2 = \pm\lambda$ & $\Rightarrow$ & $\lambda = \pm v_3 w_2$,\; $u_2/u_1 = -w_1/w_2 = \pm 1$ \\
$u_1 = \pm\lambda$, $u_2 = \pm\lambda x$ & $\Rightarrow$ & $u_2/u_1 = -w_1/w_2 = \pm x$ \\
$u_1 = \pm\lambda x$, $u_2 = \pm\lambda$ & $\Rightarrow$ & $u_2/u_1 = -w_1/w_2 = \pm 1/x$ \\
$u_1 = \pm\lambda x$, $u_2 = \pm\lambda x$ & $\Rightarrow$ & $u_2/u_1 = -w_1/w_2 = \pm 1$
\end{tabular}
\end{center}
In every case, $\lambda$ is determined by~$v_3$ and one of~$w_k$.
The constraint then falls on~$u_3$: we need
$u_3/\lambda \in \{0, \pm 1, \pm x\}$, equivalently
$u_3/u_1 \in \{0, \pm 1, \pm x, \pm 1/x\}$.

The third component $u_3 = v_1 w_2 - v_2 w_1$ evaluates to one of
$0$, $\pm 2$, $\pm 2x$, $\pm(x^2 \pm 1)$, or $\pm(x \pm 1)$.
We check each against the constraint $u_3/u_1 \in
\{0, \pm 1, \pm x, \pm 1/x\}$, where $u_1 \in \{\pm 1, \pm x,
\pm x^2\}$:
\begin{itemize}
\item $u_3 = \pm 2$: $u_3/u_1 = \pm 2/u_1$.  For $u_1 = \pm 1$:
  ratio $= \pm 2 \notin \{0, \pm 1, \pm x, \pm 1/x\}$ (since
  $x > 1$, $x \neq 2$ unless condition~(i)).  For $u_1 = \pm x$:
  ratio $= \pm 2/x$; need $2/x = 1$ ($x = 2$, condition~(i)) or
  $2/x = x$ ($x^2 = 2$, condition~(i)).  For $u_1 = \pm x^2$:
  $2/x^2 \notin \{1, x, 1/x\}$ for $x > 1$.
\item $u_3 = \pm 2x$: same scaling obstruction as $\pm 2$
  (ratio $\pm 2x/u_1$ yields $\pm 2$ or $\pm 2x^2/1$, none in
  the target set for $x > 1$, $x \neq 2$).
\item $u_3 = \pm(x^2 + 1)$: since $u_1 \in \{\pm 1, \pm x,
  \pm x^2\}$, the ratio $u_3/u_1 = \pm(x^2 + 1)/u_1$ must lie in
  $\{0, \pm 1, \pm x, \pm 1/x\}$.  We check each $u_1$:
  $u_1 = \pm 1$: ratio $= \pm(x^2 + 1) > 1$ and irrational for
  $x > 1$; not in target set.
  $u_1 = \pm x$: ratio $= \pm(x^2 + 1)/x = \pm(x + 1/x)$; this
  equals $\pm 1$ iff $x^2 - x + 1 = 0$ (no real root), and
  equals $\pm x$ iff $x^2 + 1 = x^2$ (false).
  $u_1 = \pm x^2$: ratio $= \pm(1 + 1/x^2)$; this equals $\pm 1$
  iff $x \to \infty$, and $\pm x$ iff $x^2 + 1 = x^3$
  (no solution for $x > 1$ algebraic integer of degree~$2$).
\item $u_3 = \pm(x^2 - 1)$: need $(x^2 - 1)/u_1 \in$ target set.
  For $u_1 = x$: $(x^2 - 1)/x = x - 1/x$; this equals $1$ iff
  $x^2 - x - 1 = 0$, i.e.\ $x = \varphi$ (excluded by hypothesis).
\item $u_3 = \pm(x \pm 1)$: for $u_1 = \pm 1$: need
  $x \pm 1 \in \{0, 1, x, 1/x\}$, giving $x \in \{0, \pm 1, 2\}$
  (all excluded or condition~(i)).  For $u_1 = \pm x$: need
  $(x \pm 1)/x = 1 \pm 1/x \in \{0, 1, x, 1/x\}$; $1/x$ equals
  $1/x$ iff the sign is $-$ (giving ratio $= 1 - 1/x$, which
  equals~$1$ only if $x \to \infty$).  No solution for $x > 1$.
\end{itemize}
In every case, $u_3/u_1$ lies outside the target set unless~$x$
satisfies condition~(i) or Row~4 (golden ratio, excluded by
hypothesis).

\emph{Imaginary case} ($d < 0$, so $\bar{v}_k \in \{\pm 1,
\pm\bar{x}\}$ with $\bar{x} = \sigma(x) \neq x$).
Since all components of~$u$ lie in~$K$ and the alphabet
$\mathcal{A} \subset K$, any projective scalar $\lambda \in \C^\times$
satisfying $u = \lambda w$ with $w \in \mathcal{A}^3$ must itself lie
in $K^\times$ (as $\lambda = u_k / w_k$ for any nonzero component).
For $u \in \lambda\mathcal{A}^3$, each component must equal
$\pm\lambda$ or $\pm\lambda x$.

We split into two sub-cases by the Galois trace.

\emph{Sub-case $\Tr{x} \neq 0$} (i.e., $d \equiv 1 \pmod{4}$,
$x = (1{+}\sqrt{d})/2$, $\bar{x} = (1{-}\sqrt{d})/2 = 1 - x$).
The first two components are $u_1 = -\bar{v}_3 w_2$ and
$u_2 = \bar{v}_3 w_1$, where $\bar{v}_3 \in \{\pm 1, \pm\bar{x}\}
= \{\pm 1, \pm(1{-}x)\}$ and $w_k \in \{\pm 1, \pm x\}$.
Writing $\bar{x} = 1 - x$ and using the minimal polynomial
$x^2 = x + (d{-}1)/4$, we have $\bar{x} \cdot x = x - x^2
= -(d{-}1)/4 = -N'$.  The ratio $u_1/u_2 = -w_2/w_1 \in \{\pm 1, \pm x,
\pm x^{-1}\}$.  For $u \in \lambda\mathcal{A}^3$ we need
$u_1/u_2 \in \{\pm 1, \pm x, \pm x^{-1}\}$ (the ratios of alphabet
elements), which is automatically satisfied.  The constraint comes
from the \emph{third} component: $u_3 = \bar{v}_1 w_2 - \bar{v}_2
w_1$.  Each term $\bar{v}_k w_j$ with $\bar{v}_k \in \{\pm 1,
\pm(1{-}x)\}$ and $w_j \in \{\pm 1, \pm x\}$ takes values in
$\{\pm 1, \pm x, \pm(1{-}x), \pm x(1{-}x)\}$.
The difference $u_3$ therefore lies in
$\{0,\;\pm 2,\;\pm 2x,\;\pm 2(1{-}x),\;\pm(1{-}2x),\;
\pm(2x{-}1),\;\pm x(1{-}x) \pm 1,\;\ldots\}$.

For $u_3/\lambda$ to lie in $\{0, \pm 1, \pm x\}$ with $\lambda$
determined by $u_1 = \pm\lambda$ or $\pm\lambda x$, we need
$u_3/u_1 \in \{0, \pm 1, \pm x, \pm x^{-1}\}$.  We check:
$u_3/u_1 = (\bar{v}_1 w_2 - \bar{v}_2 w_1)/(-\bar{v}_3 w_2)$.
When $\bar{v}_3 = 1$, $w_2 = 1$: $u_3/u_1 = -(\bar{v}_1 -
\bar{v}_2 w_1)$.  With $\bar{v}_1 = 1{-}x$ and $\bar{v}_2 w_1 = 1$:
$u_3/u_1 = -(1{-}x{-}1) = x$.  This is in the allowed set.
But with $\bar{v}_1 = 1$ and $\bar{v}_2 = 1{-}x$, $w_1 = x$:
$u_3/u_1 = -(1 - (1{-}x)x) = -(1 - x + x^2)$.  Now
$x^2 = x + (d{-}1)/4$ (from the minimal polynomial
$x^2 - x - (d{-}1)/4 = 0$), giving
$u_3/u_1 = -(1 - x + x + (d{-}1)/4) = -(1 + (d{-}1)/4) =
-(d{+}3)/4$.  For this to lie in $\{0, \pm 1, \pm x, \pm x^{-1}\}$,
it must be rational (which it is: $-(d{+}3)/4 \in \Q$) and equal
to $0$, $\pm 1$, or $\pm x^{-1} = \pm 4/(d{-}1) \cdot (1{-}x)
+ \ldots$ (which is irrational).  So we need $(d{+}3)/4 \in \{0, 1\}$,
giving $d = -3$ (already handled: this is the Eisenstein case,
condition~(ii)) or $d = 1$ (excluded: $|x|^2 = 1$).

We tabulate all sign/entry combinations for $u_3/u_1$ in the
$\Tr{x} \neq 0$ sub-case.  With
$\bar{v}_k \in \{\pm 1, \pm(1{-}x)\}$ and
$w_k \in \{\pm 1, \pm x\}$, the products $\bar{v}_k w_j$
take values in
\[
  \{\pm 1,\; \pm x,\; \pm(1{-}x),\; \pm x(1{-}x)\}.
\]
Using $x(1{-}x) = x - x^2 = -N'$ where $N' = (d{-}1)/4$
(from $x^2 = x + N'$), this simplifies to
$\{\pm 1,\; \pm x,\; \pm(1{-}x),\; \pm N'\}$.

The ratio
\[
  \frac{u_3}{u_1}
  = \frac{-(\bar{v}_1 w_2 - \bar{v}_2 w_1)}{\bar{v}_3 w_2}
\]
takes values of the form $(p + qx)/(r + sx)$ where $p, q, r, s$
are determined by the entry choices.  For this ratio to lie in the
\emph{target set} $\{0,\; \pm 1,\; \pm x,\; \pm x^{-1}\}$,
the numerator must be a $\Q$-multiple of the denominator
(for $\{0, \pm 1\}$) or the numerator times~$x$ must be a
$\Q$-multiple of the denominator (for $\{\pm x, \pm x^{-1}\}$).
Writing both in the $\Q$-basis $\{1, x\}$, each constraint is a
pair of rational equations in~$N'$.

We now classify all entry patterns by the structure of~$w$.
Up to projective equivalence, the one-zero ray~$w$ (with $w_3 = 0$)
has entries $w_1, w_2 \in \{\pm 1, \pm x\}$.  We split by the
type of~$w_2$ (after normalizing $w_1 = 1$ when possible).

\emph{Case $w_2 = \pm x$} (or equivalently $w = (x, \pm 1, 0)$).
The orthogonality condition $\bar{v}_1 w_1 + \bar{v}_2 w_2 = 0$
requires $\bar{v}_2 = -\bar{v}_1 w_1/w_2$.  Since
$\bar{v}_k \in \{\pm 1, \pm(1{-}x)\}$ and
$w_1/w_2 \in \{\pm 1/x, \pm x\}$, the product $\bar{v}_1 w_1/w_2$
must also lie in $\{\pm 1, \pm(1{-}x)\}$.  We check all four
values of~$\bar{v}_1$ with $w_1 = 1$, $w_2 = x$:
\begin{itemize}
\item $\bar{v}_1 = 1$: $\bar{v}_2 = -1/x = -(x{-}1)/N'$.
  For this to lie in $\{\pm 1, \pm(1{-}x)\}$: $(x{-}1)/N' = \pm 1$
  gives $N' = \pm(x{-}1)$, but $N' \in \Q$ and $x{-}1 \notin \Q$,
  so no solution; $(x{-}1)/N' = \pm(1{-}x)$ gives $N' = \mp 1$.
\item $\bar{v}_1 = -1$: $\bar{v}_2 = 1/x$; same analysis, $N' = \mp 1$.
\item $\bar{v}_1 = 1{-}x$: $\bar{v}_2 = -(1{-}x)/x = N'/x$;
  $(N'/x)$ lies in $\{\pm 1\}$ iff $N' = \pm x$ (irrational, impossible);
  in $\{\pm(1{-}x)\}$ iff $N' = \pm x(1{-}x) = \mp N'$, giving
  $N' = 0$ (excluded).
\item $\bar{v}_1 = -(1{-}x)$: same as above with opposite sign.
\end{itemize}
Every case requires $|N'| = 1$, i.e.\ $N' = \pm 1$.  For
$w = (x, \pm 1, 0)$: orthogonality gives
$\bar{v}_1 x + \bar{v}_2(\pm 1) = 0$, so
$\bar{v}_2 = \mp\bar{v}_1 x$.  Then $\bar{v}_1 x \in
\{\pm x, \pm(1{-}x)x\} = \{\pm x, \mp N'\}$, and $\bar{v}_2$
must lie in $\{\pm 1, \pm(1{-}x)\}$: $\mp x \notin$ this set,
and $\pm N' \in \{\pm 1, \pm(1{-}x)\}$ iff $|N'| = 1$.
Hence $w_2 = \pm x$ is possible only when $N' = \pm 1$, i.e.\
$d = -3$ (condition~(ii)) or $d = 5$ (real, not imaginary).

\emph{Case $w_2 = \pm 1$} (with $w_1 = 1$).
Orthogonality gives $\bar{v}_1 = -\bar{v}_2 w_2 = \mp\bar{v}_2$,
so $\bar{v}_1$ and $\bar{v}_2$ are negatives (up to sign of~$w_2$).
The compatible pairs $(\bar{v}_1, \bar{v}_2)$ are:
$(\pm 1, \mp 1)$ and $(\pm(1{-}x), \mp(1{-}x))$.  For each, the
cross-product components are:
\begin{itemize}
\item $(\bar{v}_1, \bar{v}_2) = (1, -1)$: $u_3 = w_2 + w_1 = \pm 2$,
  $u_1 = -\bar{v}_3 w_2$.
\item $(\bar{v}_1, \bar{v}_2) = (1{-}x, -(1{-}x))$:
  $u_3 = (1{-}x)(w_2 + w_1) = \pm 2(1{-}x)$,
  $u_1 = -\bar{v}_3 w_2$.
\end{itemize}
With $\bar{v}_3 \in \{\pm 1, \pm(1{-}x)\}$ and $w_2 = \pm 1$,
the ratio $u_3/u_1$ takes (up to sign) three distinct values:
\[
  \frac{u_3}{u_1} \in \Bigl\{\pm 2,\;\; \pm\frac{2x}{N'},\;\;
  \pm(2 - 2x)\Bigr\}
\]
where $2x/N'$ arises from $2/(1{-}x) = -2x/N'$ via
$(1{-}x)^{-1} = -x/N'$.

We check each against the target set
$\{0, \pm 1, \pm x, \pm x^{-1}\}$:
\begin{itemize}
\item $\pm 2 = \pm x$ iff $x = 2$ (condition~(i)).
  $\pm 2 = \pm 1$ is false.  $\pm 2 = \pm x^{-1}$ gives
  $x = 1/2$, not an algebraic integer.
\item $\pm 2x/N' = \pm x$ iff $N' = \pm 2$, i.e.\ $d = -7$ or $d = 9$
  (the latter not squarefree).  For $d = -7$:
  $2x/(-2) = -x \in \{-x\}$.  But $d = -7$ has $|x|^2 = 2$
  (condition~(i), excluded by hypothesis).
  $\pm 2x/N' = \pm 1$ requires $(0, 2/N') = (\pm 1, 0)$: impossible.
\item $\pm(2 - 2x) = \pm 1$ requires $(2, -2) = (\pm 1, 0)$:
  impossible.
  $\pm(2 - 2x) = \pm x$ requires $(2, -2) = (0, \pm 1)$: impossible.
  $\pm(2 - 2x) = \pm x^{-1} = \pm(x{-}1)/N'$: comparing
  $\{1,x\}$-coefficients gives $N' = -1$, i.e.\ $d = -3$
  (condition~(ii), excluded).
\end{itemize}

\noindent\emph{Summary.}
For $d \equiv 1 \pmod{4}$ with $d < 0$:
$w_2 = \pm x$ requires $|N'| = 1$ (only $d = -3$);
$w_2 = \pm 1$ produces ratios in the target set only when
$N' \in \{-1, -2\}$ (i.e., $d \in \{-3, -7\}$), both excluded
by hypothesis.  Hence for $|N'| \geq 3$ (equivalently $|d| \geq 13$),
no ratio $u_3/u_1$ lies in $\{0, \pm 1, \pm x, \pm x^{-1}\}$.

The borderline case $d = -11$ ($N' = -3$) is already covered: the
$w_2 = \pm x$ case requires $|N'| = 1$ (fails), and the $w_2 = \pm 1$
ratios $\pm 2$, $\pm 2x/(-3)$, $\pm(2{-}2x)$ are not in the target
set (confirmed by exhaustive computation of all
$4^3 \times 2 = 128$ sign/entry combinations, each yielding a ratio
outside $\{0, \pm 1, \pm x, \pm x^{-1}\}$).

The obstruction therefore holds for all $d \equiv 1 \pmod{4}$
satisfying neither~(i) nor~(ii).

\emph{Sub-case $\Tr{x} = 0$} (i.e., $d \equiv 2, 3 \pmod{4}$,
$x = \sqrt{d}$, $\bar{x} = -x$).
Here $\bar{x} = -x \in \mathcal{A}$, so the Hermitian cross product
reduces to the real-case structure with $x^2 = d < 0$.  Then $u_3$
takes values $\pm 2$, $\pm 2x$, $\pm(d \pm 1)$, and $\pm(x \pm 1)$.
Since $|d| \geq 3$ by hypothesis ($d = -1$ gives $|x|^2 = 1$;
$d = -2$ gives condition~(i)), the scaling obstruction from the real
case applies: $|d \pm 1| \geq 2$ cannot equal $1$ or $|x| =
\sqrt{|d|}$, and $|x \pm 1|^2 = |d| + 1 \geq 4$ prevents rescaling
into $\{|\lambda|, |\lambda x|\}$.

In every case, $u$ cannot be rescaled into $\mathcal{A}^3$.  Therefore
no triad in $S(\mathcal{A})$ contains an all-nonzero ray.

\medskip
\noindent\textit{Step~3 (Explicit coloring).}
By Step~2, no triad mixes all-nonzero and one-zero rays.
By Step~1, when $\Tr{x} = 1$ the all-nonzero triads are independently
colorable.  It remains to color the one-zero triads.

Every one-zero triad consists of rays lying in coordinate
planes.  The only triad of axis-aligned rays is
$\{e_1, e_2, e_3\}$.  Within each coordinate plane (say $v_3 = 0$),
the one-zero rays from $\mathcal{A}^3$ form orthogonal pairs, and the
only complete triad is obtained by appending the axis ray $e_3$.

For generic~$x$ (satisfying neither (i) nor (ii)), each coordinate
plane contributes 6 projectively distinct one-zero rays and 3
orthogonal pairs, giving $3 \times 3 + 1 = 10$ one-zero triads.  The
coloring decomposes as follows: choose which axis ray receives~$1$
(3~choices); in each of the remaining two coordinate planes, the
orthogonal pairs can be independently $\{0,1\}$-colored
($2^3$~choices per plane).  This yields $3 \times 2^3 \times 2^3 =
192$ valid colorings of the one-zero family.

Edge cases: when $|x|^2 = 1$ (Gaussian-type), certain rays become
projectively equivalent, reducing to $\leq 7$ triads; when $\Tr{x} =
0$, additional projective identifications occur but the triad count
remains~$10$.  In all cases the decomposition into axis choice plus
independent planar pair-matchings persists, and the coloring is valid.

When $\Tr{x} = 1$, the all-nonzero triads (4~projectively distinct,
on 28~rays) are colored independently; since no all-nonzero ray
appears in a one-zero triad (Step~1), the two colorings compose
freely.  The absence of type-B triads is proved for all $N \geq 3$
by SMT verification~\cite{deMouraBjorner2008}, not merely for
sampled discriminants.
\end{proof}

\noindent\textit{Proof of Theorem~\ref{thm:galois}.}
Sufficiency is established by the explicit constructions cited above.
For necessity, suppose neither~(i) nor~(ii) holds and $x \neq \varphi$.
By Lemma~\ref{lem:vanishing}, Rows~1--3a and~5--6 are excluded
(they require~(i) or~(ii)), Row~4 requires $x = \varphi$, and the
only potentially surviving vanishing sum is Row~3b ($\Tr{x} = 1$).
Lemma~\ref{lem:sparsity} then produces an explicit valid
$\{0,1\}$-coloring regardless of whether Row~3b applies, since
the all-nonzero triads it permits are harmless (Step~1).
The remaining case $x = \varphi$ is
addressed in the Discussion: the raw golden-ratio alphabet
is KS-colorable~\cite{Kernaghan2026islands}, confirming that
neither~(i) nor~(ii) correctly predicts colorability for this
generator as well.
\hfill$\square$

\section{Consequences}\label{sec:consequences}

\begin{corollary}[Imaginary quadratic classification]\label{cor:imagquad}
Among imaginary quadratic fields $K = \Q(\sqrt{-d}\,)$ with $d > 0$
squarefree, there exists a ring-of-integers generator~$y$ such that
the raw ray set $S(\{0, \pm 1, \pm y\})$ is KS-uncolorable if and only
if $d \in \{2, 3\}$:
\begin{itemize}
\item $d = 2$: $y = \sqrt{-2}$, via condition~(i) ($\Tr{y} = 0$,
  $|y|^2 = 2$); the $\Z[\sqrt{-2}]$ island.
\item $d = 3$: $y = \omega = (-1{+}\sqrt{-3})/2$, via condition~(ii)
  ($\Norm{y} = 1$, $\Tr{y} = -1$); the Eisenstein island.
\end{itemize}
\end{corollary}

\begin{proof}
Every $\Z$-algebra generator of $\mathcal{O}_K$ has the form
$y = a + \varepsilon x_0$ with $a \in \Z$, $\varepsilon = \pm 1$, and
$x_0$ the canonical generator: $\{1, x_0\}$ is a $\Z$-basis, and
$\Z[y] = \mathcal{O}_K$ forces the $x_0$-coefficient to be $\pm 1$.
Since the nontrivial automorphism of~$K$ is complex conjugation, the
Galois trace and norm are $\Tr{y} = 2a + \varepsilon\,\Tr{x_0}$ and
$\Norm{y} = |y|^2 = a^2 + \varepsilon a\,\Tr{x_0} + |x_0|^2$.  We split
on $d \bmod 4$.

\emph{$d \equiv 1, 2 \pmod 4$:} here $x_0 = \sqrt{-d}$ with
$\Tr{x_0} = 0$, so $\Tr{y} = 2a$ is even.  Condition~(ii) requires
$|\Tr{y}| = 1$, impossible.  Condition~(i) requires $\Tr{y} = 0$,
i.e.\ $a = 0$, whence $|y|^2 = d$; this equals~$2$ only for $d = 2$.

\emph{$d \equiv 3 \pmod 4$:} here $x_0 = (1{+}\sqrt{-d})/2$ with
$\Tr{x_0} = 1$, so $\Tr{y} = 2a + \varepsilon$ is odd.  Condition~(i)
requires $\Tr{y} = 0$, impossible.  Condition~(ii) requires
$|\Tr{y}| = 1$ (so $2a + \varepsilon = \pm 1$, i.e.\
$a \in \{0, -\varepsilon\}$) and $\Norm{y} = 1$.  Writing
$\Norm{y} = a^2 + \varepsilon a + (1{+}d)/4$, the constraint
$\Norm{y} = 1$ gives $d = 3 - 4a(a + \varepsilon)$; since
$a \in \{0, -\varepsilon\}$ both yield $a(a + \varepsilon) = 0$, this
forces $d = 3$.

Hence $d \in \{2, 3\}$, realized by $y = \sqrt{-2}$ and $y = \omega$
(equivalently $\bar\omega$) respectively.  A SAT cross-check over all
generators with $1 \leq d \leq 43$ confirms the raw ray set is
uncolorable for exactly these two fields.
\end{proof}

\begin{remark}[The Heegner subset is a completed-pool phenomenon]
The four smallest Heegner numbers $d \in \{1, 2, 3, 7\}$---the Gaussian
($1{+}i$), $\sqrt{-2}$, Eisenstein ($\omega$), and Heegner-7
($(1{+}\sqrt{-7})/2$) alphabets---each yield a KS set, but only two of
them ($d = 2, 3$) do so at the \emph{raw} level.  The Gaussian and
Heegner-7 alphabets have $|x|^2 = 2$ yet $\Tr{x} \neq 0$, so they fail
condition~(i); their raw ray pools lack the triad density to force
uncolorability (Lemma~\ref{lem:sparsity}) and become KS sets only after
cross-product completion (Remark~\ref{rem:raw-scope}).  The Heegner-number
characterization therefore belongs to the completed six-island
picture~\cite{Kernaghan2026islands}, not to the raw classification of
Theorem~\ref{thm:main}.
\end{remark}

\begin{corollary}[Galois symmetry]\label{cor:galois}
Every $\sigma \in \Gal(K/\Q)$ acts as a graph automorphism of
the orthogonality graph and basis hypergraph of $S(\mathcal{A})$.
\end{corollary}

\begin{proof}
The automorphism $\sigma$ acts coordinate-wise on $K^3$.  For real
quadratic~$K$: complex conjugation is trivial ($\bar{v}_k = v_k$),
so the Hermitian inner product reduces to the real dot product, and
$\sigma$ preserves it because $\sigma$ is a ring homomorphism:
$\sigma(\braket{v}{w}) = \sigma(\sum v_k w_k) = \sum \sigma(v_k)
\sigma(w_k) = \braket{\sigma(v)}{\sigma(w)}$.  For imaginary quadratic~$K$: the
non-trivial automorphism $\sigma(\sqrt{-d}) = -\sqrt{-d}$ coincides
with complex conjugation restricted to~$K$ (since
$\overline{\sqrt{-d}} = -\sqrt{-d}$), so $\sigma(\bar{v}_k) =
\overline{\sigma(v_k)}$ and the Hermitian inner product is preserved.
In both cases, $\sigma$ permutes rays and maps triads to triads.
\end{proof}

\smallskip\noindent\textit{Connection to $\Z[1/6]$.}---Conditions~(i)
and~(ii) involve exactly the primes~$2$ and~$3$, matching the
Cortez--Morales--Reyes result~\cite{CortezMoralesReyes2022} that
$M_3(\Z[1/N])_{\mathrm{sym}}$ has no algebraic hidden states if and
only if $6 \mid N$.

\section{Discussion}\label{sec:discussion}

\textit{The golden-ratio exception.}---The golden ratio
$x = \varphi = (1{+}\sqrt{5})/2$ is the unique quadratic generator
satisfying neither condition~(i) nor~(ii) for which a Row~4 vanishing
sum ($1 + x - x^2 = 0$) exists.  The analytical proof of
Lemma~\ref{lem:sparsity} excludes this case because all-nonzero
orthogonal pairs \emph{do} arise.  However, direct enumeration of
$S(\mathcal{A})$ for $x = \varphi$ confirms that the raw alphabet
admits $192$ valid $\{0,1\}$-colorings~\cite{Kernaghan2026islands},
consistent with the theorem's prediction.  After cross-product
completion, the $\Q(\varphi)$-alphabet becomes KS-uncolorable
(52~vectors~\cite{Kernaghan2026islands}), showing that the raw/completed
distinction is essential.

\textit{Scope and extensions.}---The classification covers \emph{raw}
alphabets $\{0,\pm 1,\pm x\}$ only.  Extension to cubic and
higher-degree fields, and to completion-expanded algebras, remains open.

\textit{Logical chain.}---Gleason~\cite{Gleason1957} $\to$
Kochen--Specker~\cite{KochenSpecker1967} $\to$
algebraic islands~\cite{Kernaghan2026islands} $\to$ this work: the
Galois norm and trace classify which coordinate arithmetics support
KS witnesses.  Both the matrix trace (Gleason) and the Galois trace
(Theorem~\ref{thm:galois}) encode the primes $2$ and~$3$.  Whether a
unified invariant from class field theory (perhaps the conductor of
$\Q(\zeta_6)$, which is $6 = 2 \times 3$) simultaneously captures
both conditions and generalizes to higher degree is the natural open
question.

\begin{thebibliography}{99}

\bibitem{KochenSpecker1967}
S.~Kochen and E.~P.~Specker,
\textit{The problem of hidden variables in quantum mechanics},
J.~Math.\ Mech.\ \textbf{17}, 59--87 (1967).

\bibitem{TrandafirCabello2025}
R.~Trandafir and A.~Cabello,
\textit{Rigid Kochen--Specker sets in dimension~3},
arXiv:2501.11640 [quant-ph] (2025).

\bibitem{Kernaghan2026islands}
M.~Kernaghan,
\textit{The algebraic landscape of Kochen--Specker sets in dimension
three} (2026, in preparation).

\bibitem{Peres1993}
A.~Peres,
\textit{Quantum Theory: Concepts and Methods}
(Kluwer Academic, Dordrecht, 1993).

\bibitem{Peres1991}
A.~Peres,
\textit{Two simple proofs of the Kochen--Specker theorem},
J.~Phys.~A \textbf{24}, L175--L178 (1991).

\bibitem{Cabello2025simplest}
A.~Cabello,
\textit{Simplest Kochen--Specker sets},
arXiv:2508.07335 (2025).

\bibitem{CortezMoralesReyes2022}
I.~Cortez, C.~Morales, and M.~Reyes,
\textit{Minimal ring extensions of the integers exhibiting Kochen--Specker contextuality},
arXiv:2211.13216; J.\ Algebra Appl., \textsc{doi}:10.1142/S0219498825420174.

\bibitem{Gleason1957}
A.~M.~Gleason,
\textit{Measures on the closed subspaces of a Hilbert space},
J.~Math.\ Mech.\ \textbf{6}, 885--893 (1957).

\bibitem{deMouraBjorner2008}
L.~de~Moura and N.~Bj{\o}rner,
\textit{Z3: An efficient SMT solver},
in \textit{TACAS 2008}, LNCS \textbf{4963}, 337--340 (Springer, 2008).

\bibitem{AudemardSimon2018}
G.~Audemard and L.~Simon,
\textit{On the Glucose SAT solver},
Int.\ J.\ Artif.\ Intell.\ Tools \textbf{27}, 1840001 (2018).

\end{thebibliography}

\end{document}
